\subsection{Формула шлифовщика}
\label{sec:grinder-formula}

\begin{wrapfigure}[11]{r}{0.5\tw}
    \centering
    \vspace{-1pc}
    \tikzsetnextfilename{mirror-sagitta}
    \begin{tikzpicture}
        \footnotesize

        \draw [dashes] (-3.9, 0) -- (0.7, 0);
        \draw [semithick] (-3, 0) ++(-30:3) arc(-30:30:3);
        \draw [dashes] (-0.402, -1.5) -- (-0.402, 1.5);
        \draw (-3, 0) -- (-0.402, 1.5);

        \draw [latex-latex] (-0.402, 0) -- (0, 0);
        \draw [latex-latex] (-1.1, -1.5) -- (-1.1, 1.5);
        \draw (-1.3, -1.5) -- (-0.402, -1.5);
        \draw (-1.3, 1.5) -- (-0.402, 1.5);

        \draw (-0.2, 0) node[anchor=south] {$h$};
        \draw (-1.1, 0.4) node[anchor=east] {$D$};
        \draw (-1.9, 0.75) node[anchor=south east] {$R$};
        \draw (-3, 0) node[anchor=north] {$C$};
        \draw (0, 0) node[anchor=north west] {$O$};

        \draw [fill=white] (-3, 0) circle (.03);
        \draw [fill=white] (0, 0) circle (.03);
    \end{tikzpicture}
    \caption{Сечение сферического зеркала диаметром $D$ и радиусом кривизны $R$; $h$~--- стрелка прогиба}
    \label{pic:mirror-sagitta}
\end{wrapfigure}
Главное зеркало рефлектора изготавливают шлифовкой: на стеклянном диске диаметром $D$ вышлифовывают вогнутую сферическую поверхность радиуса кривизны $R$, которую затем при необходимости доводят до параболоида. Контролировать форму поверхности в процессе работы удобно по \term{стрелке прогиба} (стрелке кривизны) $h$~--- глубине выемки в центре зеркала относительно его края, \lookPicRef{pic:mirror-sagitta}. Соотношение, связывающее стрелку прогиба с диаметром и радиусом кривизны зеркала, называют \term{формулой шлифовщика}. Подробнее о технологии изготовления зеркал см.~\cite{sikoruk}.

Рассмотрим сечение зеркала плоскостью, проходящей через оптическую ось. Оно представляет собой дугу окружности радиуса $R$ с центром в точке $C$, стягиваемую хордой длины $D$. Расстояние от центра кривизны до хорды равно $R - h$, поэтому по теореме Пифагора
\begin{gather*}
    R^2 = \left(\frac{D}{2}\right)^2 + (R - h)^2,\\
    h = R - \sqrt{R^2 - \frac{D^2}{4}}.
\end{gather*}
Полученное выражение точное, однако на практике всегда $D \ll R$, и его удобно упростить. Разложим корень в ряд по малому параметру $D/2R$, ограничившись первым отличным от нуля слагаемым:
\begin{equation}
    h = R \left(1 - \sqrt{1 - \frac{D^2}{4R^2}}\right)
      \approx R \cdot \frac{D^2}{8R^2}
      = \frac{D^2}{8R}.
    \label{eq:grinder-formula}
\end{equation}
Фокус сферического зеркала лежит посередине между зеркалом и центром его кривизны, $F = R/2$, \lookPicRef{pic:sphere-aberrations-mirrow}. Следовательно, формулу шлифовщика можно записать через фокусное расстояние или относительное отверстие $\forall = D/F$:
\begin{equation}
    h \approx \frac{D^2}{16F} = \frac{D \forall}{16}.
\end{equation}
Например, для зеркала диаметром $150$~мм с фокусным расстоянием $1200$~мм ($\forall = 1/8$) стрелка прогиба составляет всего $h \approx 1.2$~мм. Объём стекла, снимаемого при шлифовке, равен объёму шарового сегмента высоты $h$, то есть
\begin{equation*}
    V = \frac{\pi h^2 (3R - h)}{3} \approx \pi h^2 R \approx \frac{\pi D^2 h}{8} \approx 10\text{ см}^3.
\end{equation*}

Обратная задача возникает при измерении радиуса кривизны готовой поверхности \term{сферометром}~--- прибором из трёх опорных ножек, расположенных на окружности диаметра $d$, и измерительного стержня в её центре. Сферометр ставят на зеркало и измеряют стрелку прогиба $h$ на диаметре $d$, откуда, повторяя выкладки выше без приближений,
\begin{equation}
    R = \frac{d^2}{8h} + \frac{h}{2}.
\end{equation}
Второе слагаемое~--- поправка к формуле шлифовщика \eqref{eq:grinder-formula}, которой при $h \ll d$ можно пренебречь.

\paragraph{Сфера или параболоид?} Параболическое зеркало свободно от сферической аберрации, \lookSecRef{sec:aberrations}, но параболизация~--- самый трудоёмкий этап изготовления. Выясним, когда сферическое зеркало можно использовать вместо параболического. Пусть $r$~--- расстояние от оптической оси, $z$~--- отклонение поверхности от касательной плоскости в вершине зеркала. Уравнение параболы с радиусом кривизны $R$ в вершине имеет вид $z_\text{п} = r^2/2R$. Разложим в ряд уравнение сферы того же радиуса, удержав на этот раз два слагаемых:
\begin{equation*}
    z_\text{сф} = R - \sqrt{R^2 - r^2}
        = \frac{r^2}{2R} + \frac{r^4}{8R^3} + \dots
\end{equation*}
Таким образом, сфера отклоняется от параболоида на величину $\Delta z = r^4/8R^3$, максимальную на краю зеркала, при $r = D/2$:
\begin{equation*}
    \Delta z_\text{max} = \frac{D^4}{128 R^3} = \frac{D^4}{1024 F^3}.
\end{equation*}
При отражении разность хода лучей удваивается. С другой стороны, подбором радиуса сферы, что равносильно перефокусировке, размах ошибки волнового фронта можно уменьшить ещё в четыре раза.\footnote{Ошибка волнового фронта сферы относительно параболоида $W(r) = a r^4 - \beta r^2$, где $a = 1/4R^3$, а слагаемое $\beta r^2$ описывает изменение радиуса сферы. Размах $W$ на отрезке $[0, D/2]$ минимален при $\beta = a D^2/4$: тогда $W(0) = W(D/2) = 0$, а минимум $W = -a D^4/64$ достигается при $r^2 = D^2/8$. Размах $a D^4/64$ ровно в четыре раза меньше, чем $W(D/2) = a D^4/16$ при $\beta = 0$.} Наименьшая достижимая ошибка волнового фронта
\begin{equation}
    \Delta_\text{min} = \frac{2 \Delta z_\text{max}}{4} = \frac{D^4}{2048 F^3}.
\end{equation}
Согласно критерию Рэлея оптическая система строит практически идеальное изображение, если ошибка волнового фронта не превышает $\lambda/4$. Отсюда условие, при котором сферическое зеркало не требует параболизации:
\begin{equation}
    \frac{D^4}{2048 F^3} \leqslant \frac{\lambda}{4}
    \quad \Longrightarrow \quad
    F \geqslant \sqrt[3]{\frac{D^4}{512 \lambda}}.
    \label{eq:sphere-vs-paraboloid}
\end{equation}
Для зеркала диаметром $150$~мм и $\lambda = 550$~нм получаем $F \geqslant 1.2$~м, то есть $\forall \leqslant 1/8$. С ростом диаметра требование ужесточается: зеркало диаметром $300$~мм остаётся сферическим лишь при $\forall \leqslant 1/10$, поэтому крупные зеркала практически всегда параболизуют.
